% i18n_qa_sv_paper.tex — Swedish physics paper on quantum mechanics
% Intentional errors: wrong quotes (English "..." instead of Swedish >>...<<),
%   decimal period instead of decimal comma (3.14 vs 3,14),
%   aa/ae/oe encoding issues
\documentclass[12pt,a4paper]{article}

\usepackage[swedish]{babel}
\usepackage[utf8]{inputenc}
\usepackage[T1]{fontenc}
\usepackage{amsmath}
\usepackage{amssymb}
\usepackage{amsthm}
\usepackage{graphicx}
\usepackage{hyperref}
\usepackage[margin=2.5cm]{geometry}
\usepackage{booktabs}
\usepackage{natbib}
\usepackage{enumitem}
\usepackage{braket}
\usepackage{microtype}
\usepackage{xcolor}
\usepackage{float}
\usepackage{siunitx}
\usepackage{physics}
\usepackage{mathrsfs}

\theoremstyle{plain}
\newtheorem{theorem}{Sats}[section]
\newtheorem{lemma}[theorem]{Lemma}
\newtheorem{proposition}[theorem]{P\r{a}st\r{a}ende}

\theoremstyle{definition}
\newtheorem{definition}[theorem]{Definition}
\newtheorem{example}[theorem]{Exempel}

\theoremstyle{remark}
\newtheorem{remark}[theorem]{Anm\"arkning}

\title{Kvantmekaniska sammanfl\"atningar och Bellska olikheter\\
i system med flera partiklar}
\author{%
  Erik Lindstr\"om\thanks{Institutionen f\"or fysik, Lunds universitet, Lund.}
  \and
  Anna S\"oderstr\"om\thanks{Institutionen f\"or teoretisk fysik,
  Kungliga Tekniska h\"ogskolan, Stockholm.}
}
\date{den 6 april 2026}

\begin{document}

\maketitle

\begin{abstract}
I denna artikel unders\"oker vi kvantmekanisk sammanfl\"atning
(entanglement) i system best\r{a}ende av tre eller fler partiklar.
Vi h\"arleder en ny familj av Bellska olikheter som \"ar striktare
\"an de klassiska CHSH-olikheterna f\"or flerpartikelsystem.
V\r{a}ra teoretiska resultat visar att det maximala br\"attet mot
de lokala realistiska gr\"anserna \"okar exponentiellt med antalet
partiklar. Vi presenterar ocks\r{a} numeriska simuleringar som
bekr\"aftar de teoretiska f\"oruts\"agelserna f\"or system med upp till
% INTENTIONAL ERROR: decimal period instead of comma
12 partiklar. Den maximala kr\"ankningskvoten \"ar 2.83 f\"or tv\r{a}
partiklar och v\"axer till 45.7 f\"or tio partiklar. Resultaten
har implikationer f\"or kvantinformationsbehandling och
kvantkryptografi.
\end{abstract}

\tableofcontents

%% ---------------------------------------------------------------
\section{Inledning}
\label{sec:inledning}

Kvantmekanisk sammanfl\"atning \"ar ett av de mest fascinerande
fenomenen i modern fysik. Sedan Einstein, Podolsky och Rosens
(EPR) ber\"omda artikel 1935~\citep{einstein1935} har fr\r{a}gan
om kvantmekanikens fullst\"andighet debatterats intensivt.

% INTENTIONAL ERROR: English quotes instead of Swedish >>...<<
John Bells banbrytande arbete 1964~\citep{bell1964} visade att
ingen "lokal realistisk" teori kan reproducera alla
f\"oruts\"agelser fr\r{a}n kvantmekaniken. Bells olikhet ger ett
experimentellt testbart kriterium f\"or att skilja mellan
kvantmekanik och lokala dolda variabelmodeller.

F\"or tv\r{a} partiklar i ett singletttillst\r{a}nd
\begin{equation}
\label{eq:singlet}
\ket{\Psi^-} = \frac{1}{\sqrt{2}} \left( \ket{\uparrow\downarrow} - \ket{\downarrow\uparrow} \right),
\end{equation}
ger CHSH-olikheten~\citep{clauser1969}:
\begin{equation}
\label{eq:chsh}
|S| = |\langle A_1 B_1 \rangle + \langle A_1 B_2 \rangle + \langle A_2 B_1 \rangle - \langle A_2 B_2 \rangle| \leq 2,
\end{equation}
% INTENTIONAL ERROR: decimal period instead of comma
d\"ar den lokala realistiska gr\"ansen \"ar 2. Kvantmekaniken
f\"oruts\"ager ett maximalt v\"arde p\r{a} $2\sqrt{2} \approx 2.828$,
k\"ant som Tsirelsons gr\"ans\footnote{Boris Tsirelson (1980)
visade att detta \"ar den maximala kr\"ankningen m\"ojlig inom
kvantmekaniken.}.

V\r{a}rt m\r{a}l \"ar att generalisera dessa resultat till system
med $N \geq 3$ partiklar och h\"arleda nya, striktare olikheter.

\subsection{Tidigare arbeten}

Greenberger, Horne och Zeilinger (GHZ)~\citep{ghz1989} visade
att f\"or tre partiklar finns ett \"annu st\"orre gap mellan
kvantmekanik och lokal realism. GHZ-tillst\r{a}ndet
\begin{equation}
\label{eq:ghz}
\ket{\mathrm{GHZ}} = \frac{1}{\sqrt{2}} \left( \ket{000} + \ket{111} \right)
\end{equation}
ger ett perfekt motexempel utan statistiska argument\footnote{%
Till skillnad fr\r{a}n Bells olikhet, som kr\"aver statistik \"over
m\r{a}nga m\"atningar, ger GHZ ett enstaka motexempel.}.

Mermin~\citep{mermin1990} formulerade en generalisering av
CHSH-olikheten f\"or $N$ partiklar och visade att kvantmekaniken
% INTENTIONAL ERROR: decimal period
kr\"anker den med en faktor som v\"axer som $2^{(N-1)/2} \approx 1.414^{N-1}$.

\subsection{V\r{a}ra bidrag}

De huvudsakliga bidragen i denna artikel \"ar:

\begin{enumerate}[label=(\roman*)]
  \item En ny familj av Bellska olikheter f\"or $N$ partiklar
        (Sats~\ref{thm:huvudsats});
  \item Bevis f\"or att den maximala kvantmekaniska kr\"ankningen
        \"ar exponentiell i $N$ (Sats~\ref{thm:exponentiell});
  \item Numerisk verifiering f\"or $N \leq 12$
        (Avsnitt~\ref{sec:numeriska}).
\end{enumerate}

%% ---------------------------------------------------------------
\section{Teoretisk bakgrund}
\label{sec:bakgrund}

\subsection{Kvantmekaniska grundbegrepp}

% INTENTIONAL ERROR: English quotes instead of Swedish
I kvantmekaniken beskrivs tillst\r{a}ndet hos ett system av
ett "tillst\r{a}ndsvektor" $\ket{\psi}$ i ett Hilbertrum
$\mathcal{H}$. F\"or ett tv\r{a}partikelsystem \"ar
tillst\r{a}ndsrummet tensorprodukten:
\begin{equation}
\label{eq:tensor}
\mathcal{H} = \mathcal{H}_A \otimes \mathcal{H}_B.
\end{equation}

En observabel representeras av en sj\"alvkonjugerad operat\"or
$\hat{O}$ med egenv\"arden $\{o_i\}$ och egentillst\r{a}nd $\{\ket{o_i}\}$.
V\"antev\"ardet \"ar:
\begin{equation}
\label{eq:expectation}
\langle \hat{O} \rangle = \bra{\psi} \hat{O} \ket{\psi} = \sum_i o_i |\braket{o_i | \psi}|^2.
\end{equation}

\subsection{Sammanfl\"atning}

\begin{definition}
\label{def:entanglement}
Ett tv\r{a}partikeltillst\r{a}nd $\ket{\psi} \in \mathcal{H}_A \otimes \mathcal{H}_B$
\"ar sammanfl\"atat om det inte kan skrivas som en tensorprodukt
$\ket{\psi} \neq \ket{\alpha} \otimes \ket{\beta}$ f\"or n\r{a}gra
$\ket{\alpha} \in \mathcal{H}_A$, $\ket{\beta} \in \mathcal{H}_B$.
\end{definition}

F\"or att kvantifiera sammanfl\"atning anv\"ands entropi. Den
von Neumannska entropin f\"or en reducerad densitetsmatris
$\rho_A = \mathrm{Tr}_B(\ket{\psi}\bra{\psi})$ \"ar:
\begin{equation}
\label{eq:von_neumann}
S(\rho_A) = -\mathrm{Tr}(\rho_A \log_2 \rho_A) = -\sum_i \lambda_i \log_2 \lambda_i,
\end{equation}
d\"ar $\{\lambda_i\}$ \"ar egenv\"ardena till $\rho_A$.

\subsection{Bellska olikheter}

\begin{definition}
\label{def:bell_operator}
Belloperatorn f\"or CHSH-olikheten \"ar:
\[
\hat{B}_{\mathrm{CHSH}} = A_1 \otimes (B_1 + B_2) + A_2 \otimes (B_1 - B_2),
\]
d\"ar $A_i, B_j$ \"ar observabler med egenv\"arden $\pm 1$.
\end{definition}

%% ---------------------------------------------------------------
\section{Nya Bellska olikheter f\"or flerpartikelsystem}
\label{sec:nya_olikheter}

\subsection{Konstruktion av Belloperatorn}

F\"or $N$ partiklar definierar vi en generaliserad Belloperator:
\begin{equation}
\label{eq:bell_N}
\hat{B}_N = \frac{1}{2} \left[ (A_1 + A_2) \otimes \hat{B}_{N-1} + (A_1 - A_2) \otimes \hat{B}_{N-1}' \right],
\end{equation}
d\"ar $\hat{B}_{N-1}$ och $\hat{B}_{N-1}'$ \"ar Belloperatorer
f\"or $N-1$ partiklar med olika m\"atinst\"allningar.

\begin{theorem}[Huvudsats]
\label{thm:huvudsats}
F\"or varje $N \geq 2$ g\"aller den lokala realistiska olikheten:
\begin{equation}
\label{eq:lr_bound}
|\langle \hat{B}_N \rangle_{\mathrm{LR}}| \leq 1.
\end{equation}
\end{theorem}

\begin{proof}
Vi bevisar med induktion \"over $N$. Basfallet $N = 2$ \"ar
CHSH-olikheten~\eqref{eq:chsh} normaliserad med en faktor
$1/2$. F\"or induktionssteget, antag att $|\langle \hat{B}_{N-1} \rangle_{\mathrm{LR}}| \leq 1$
och $|\langle \hat{B}_{N-1}' \rangle_{\mathrm{LR}}| \leq 1$. D\r{a}:
\begin{align}
|\langle \hat{B}_N \rangle_{\mathrm{LR}}|
&\leq \frac{1}{2} |a_1 + a_2| \cdot |\langle \hat{B}_{N-1} \rangle| + \frac{1}{2} |a_1 - a_2| \cdot |\langle \hat{B}_{N-1}' \rangle| \notag \\
&\leq \frac{1}{2} (|a_1 + a_2| + |a_1 - a_2|) = 1, \label{eq:induction}
\end{align}
d\"ar sista steget f\"oljer av att $a_1, a_2 \in \{-1, +1\}$.
\end{proof}

\subsection{Kvantmekanisk kr\"ankning}

\begin{theorem}
\label{thm:exponentiell}
F\"or GHZ-tillst\r{a}ndet med $N$ partiklar g\"aller:
\begin{equation}
\label{eq:qm_violation}
\langle \hat{B}_N \rangle_{\mathrm{QM}} = 2^{(N-2)/2},
\end{equation}
vilket ger en kr\"ankningskvot
% INTENTIONAL ERROR: decimal period instead of comma
$Q_N = 2^{(N-2)/2}$ som v\"axer exponentiellt. F\"or $N = 2$
\"ar $Q_2 = 1.0$, f\"or $N = 3$ \"ar $Q_3 = 1.414$, och f\"or
$N = 10$ \"ar $Q_{10} = 22.627$.
\end{theorem}

\begin{proof}
Vi v\"aljer m\"atriktningarna optimalt. F\"or GHZ-tillst\r{a}ndet
$\ket{\mathrm{GHZ}_N} = (\ket{0}^{\otimes N} + \ket{1}^{\otimes N})/\sqrt{2}$
och m\"atriktningarna $\theta_{k,1} = 0$, $\theta_{k,2} = \pi/(2N)$
f\"or varje partikel $k$, f\r{a}r vi:
\begin{align}
\langle \hat{B}_N \rangle &= \cos^N\!\left(\frac{\pi}{2N}\right) + \binom{N}{2} \cos^{N-2}\!\left(\frac{\pi}{2N}\right) \sin^2\!\left(\frac{\pi}{2N}\right) + \cdots \notag \\
&= 2^{(N-2)/2}. \label{eq:violation_calc}
\end{align}
Det fullst\"andiga beviset anv\"ander identiteter f\"or
Tjebysjovpolynom.
\end{proof}

%% ---------------------------------------------------------------
\section{Numeriska resultat}
\label{sec:numeriska}

\subsection{Ber\"akningsmetod}

Vi ber\"aknar det maximala v\"ardet av $\langle \hat{B}_N \rangle$
\"over alla m\"ojliga tillst\r{a}nd och m\"atriktningar med hj\"alp
av semidefinit programmering (SDP)~\citep{navascues2008}.

\begin{table}[ht]
\centering
% INTENTIONAL ERROR: decimal periods instead of commas throughout
\caption{Maximal kvantmekanisk kr\"ankning av Bellska olikheten}
\label{tab:violation}
\begin{tabular}{ccccl}
\toprule
$N$ & LR-gr\"ans & QM-maximum & Kr\"ankningskvot & Tillst\r{a}nd \\
\midrule
2 & 1.000 & 1.414 & 1.414 & Singlett \\
3 & 1.000 & 2.000 & 2.000 & GHZ \\
4 & 1.000 & 2.828 & 2.828 & GHZ \\
5 & 1.000 & 4.000 & 4.000 & GHZ \\
6 & 1.000 & 5.657 & 5.657 & GHZ \\
8 & 1.000 & 11.314 & 11.314 & GHZ \\
10 & 1.000 & 22.627 & 22.627 & GHZ \\
12 & 1.000 & 45.255 & 45.255 & GHZ \\
\bottomrule
\end{tabular}
\end{table}

Tabell~\ref{tab:violation} bekr\"aftar att kr\"ankningskvoten
v\"axer exponentiellt med $N$, i \"overensst\"ammelse med
Sats~\ref{thm:exponentiell}.

\subsection{Janus-tillst\r{a}nd}

% INTENTIONAL ERROR: English quotes
F\"orutom GHZ-tillst\r{a}nd har vi ocks\r{a} unders\"okt
s\r{a} kallade "W-tillst\r{a}nd":
\begin{equation}
\label{eq:w_state}
\ket{W_N} = \frac{1}{\sqrt{N}} \left( \ket{100\cdots0} + \ket{010\cdots0} + \cdots + \ket{000\cdots1} \right).
\end{equation}

\begin{figure}[ht]
\centering
\fbox{\parbox{0.75\textwidth}{%
\centering
\textbf{Kr\"ankningskvot som funktion av $N$}\\[8pt]
\begin{tabular}{ccc}
$N$ & GHZ & W-tillst\r{a}nd \\
\hline
3 & 2.000 & 1.207 \\
5 & 4.000 & 1.512 \\
7 & 8.000 & 1.693 \\
10 & 22.627 & 1.841 \\
\end{tabular}\\[6pt]
\emph{GHZ-tillst\r{a}nd ger exponentiell kr\"ankning,}\\
\emph{W-tillst\r{a}nd ger logaritmisk.}
}}
\caption{J\"amf\"orelse av kr\"ankningskvoten f\"or GHZ- och W-tillst\r{a}nd
som funktion av antalet partiklar $N$. GHZ-tillst\r{a}ndet \"ar optimalt.}
\label{fig:violation}
\end{figure}

\subsection{Experimentella \"overv\"aganden}

% INTENTIONAL ERROR: aa encoding issue — using 'a' instead of '\r{a}'
F\"or experimentell verifiering kr\"avs h\"og fidelitet i
tillst\r{a}ndspreparation. Med nuvarande teknik kan GHZ-tillstand
med upp till $N = 14$ fotoner genereras~\citep{wang2018}.

\begin{figure}[ht]
\centering
\fbox{\parbox{0.75\textwidth}{%
\centering
% INTENTIONAL ERROR: decimal period
\textbf{Kr\"avd fidelitet f\"or signifikant kr\"ankning}\\[8pt]
\begin{tabular}{cc}
$N$ & Minsta fidelitet $F_{\min}$ \\
\hline
3 & 0.707 \\
5 & 0.841 \\
8 & 0.924 \\
10 & 0.951 \\
12 & 0.967 \\
\end{tabular}\\[6pt]
\emph{H\"ogre $N$ kr\"aver h\"ogre fidelitet.}
}}
\caption{Minsta tillst\r{a}ndsfidelitet $F_{\min}$ som kr\"avs
f\"or att observera en statistiskt signifikant (3$\sigma$)
kr\"ankning av den Bellska olikheten, som funktion av $N$.}
\label{fig:fidelity}
\end{figure}

%% ---------------------------------------------------------------
\section{Till\"ampningar}
\label{sec:tillamningar}

\subsection{Kvantkryptografi}

V\r{a}ra resultat har direkta till\"ampningar inom kvantkryptografi.
% INTENTIONAL ERROR: decimal period
Device-independent quantum key distribution (DI-QKD) anv\"ander
kr\"ankningar av Bellska olikheter f\"or att garantera s\"akerhet.
Den h\"ogre kr\"ankningskvoten f\"or flerpartikelsystem inneb\"ar
en nyckelhastighet p\r{a} $R \approx 0.85$ bitar per m\"atning
j\"amf\"ort med $R \approx 0.42$ f\"or tv\r{a}partikelsystem.

\subsection{Kvantber\"akning}

Flerpartikelsammanfl\"atning \"ar en grundl\"aggande resurs i
kvantber\"akning. V\r{a}ra olikheter kan anv\"andas f\"or att
verifiera sammanfl\"atning i kvantdatorer med m\r{a}nga
kvantbitar\footnote{Google Sycamore-processorn har 72 kvantbitar
och IBM Eagle har 127 kvantbitar (2024).}.

%% ---------------------------------------------------------------
\section{Slutsatser}
\label{sec:slutsatser}

Vi har h\"arlett en ny familj av Bellska olikheter f\"or
flerpartikelsystem (Sats~\ref{thm:huvudsats}) och visat att
den kvantmekaniska kr\"ankningen v\"axer exponentiellt med
antalet partiklar (Sats~\ref{thm:exponentiell}). Numeriska
simuleringar bekr\"aftar de teoretiska resultaten f\"or system
med upp till 12 partiklar (Tabell~\ref{tab:violation}).

Framtida arbete inkluderar experimentell verifiering med
fotoniska system och till\"ampningar inom flerparti-QKD-protokoll.

%% ---------------------------------------------------------------
\begin{thebibliography}{99}

\bibitem[Bell(1964)]{bell1964}
Bell, J.\,S. (1964).
On the Einstein--Podolsky--Rosen paradox.
\emph{Physics Physique Fizika}, 1(3), 195--200.

\bibitem[Clauser et al.(1969)]{clauser1969}
Clauser, J.\,F., Horne, M.\,A., Shimony, A. and Holt, R.\,A. (1969).
Proposed experiment to test local hidden-variable theories.
\emph{Physical Review Letters}, 23(15), 880--884.

\bibitem[Einstein, Podolsky och Rosen(1935)]{einstein1935}
Einstein, A., Podolsky, B. och Rosen, N. (1935).
Can quantum-mechanical description of physical reality be considered complete?
\emph{Physical Review}, 47(10), 777--780.

\bibitem[Greenberger, Horne och Zeilinger(1989)]{ghz1989}
Greenberger, D.\,M., Horne, M.\,A. och Zeilinger, A. (1989).
Going beyond Bell's theorem.
In \emph{Bell's Theorem, Quantum Theory, and Conceptions of the Universe},
Kluwer Academic Publishers, 69--72.

\bibitem[Mermin(1990)]{mermin1990}
Mermin, N.\,D. (1990).
Extreme quantum entanglement in a superposition of macroscopically distinct states.
\emph{Physical Review Letters}, 65(15), 1838--1840.

\bibitem[Navascu\'{e}s et al.(2008)]{navascues2008}
Navascu\'{e}s, M., Pironio, S. and Ac\'{i}n, A. (2008).
A convergent hierarchy of semidefinite programs characterizing the
set of quantum correlations.
\emph{New Journal of Physics}, 10, 073013.

\bibitem[Wang et al.(2018)]{wang2018}
Wang, X.-L., et al. (2018).
18-qubit entanglement with six photons' three degrees of freedom.
\emph{Physical Review Letters}, 120(26), 260502.

\end{thebibliography}

\end{document}
